
\section{算子 $\mathfrak p$、$\mathfrak q$ 与 $\mathfrak r$ 的构造细节}
\label{sec:7-4}

本节续接第~\ref{sec:3-8} 节。
由于本节还使用
$\mathcal M(M;\widetilde\beta)$，
我们用 $\mathcal M(L;\beta)$ 等记号
代替 $\mathcal M(\beta)$ 等。

\subsection{算子 $\mathfrak p$ 的构造细节}
\label{subsec:7-4-1}
\liftIndex{I285}%

首先证明命题~\ref{prop:3-8-27}。
证明结合第~\ref{sec:3-8} 节
及第~\ref{subsec:7-1-4} 小节的论证
（例如定理~\ref{thm:7-1-43}）。
此前讨论的是
$k+1\geq1$，
即至少有一个边界标记点时的
$\mathcal M^{\mathrm{main}}_{k+1,\ell}(L;\beta)$；
本小节则研究
$\mathcal M^{\mathrm{main}}_{0,\ell}(L;\beta)$。

正如第~\ref{subsec:3-8-3} 小节已经指出的，
此处的新问题是：若
$\beta\in H_2(M,L;\mathbb Z)$
\translatornote{源书此处将 $\beta$ 写作
$H^2(M,L;\mathbb Z)$ 中的类；
但紧接的长正合列及边界同态
均属于同调，故改用 $H_2$。
见疑误 A-0554。}
在自然长正合列
\[
H_2(M)\xrightarrow{i_*}H_2(M,L)
\xrightarrow{\partial}H_1(L)
\]
下满足 $\partial(\beta)=0$，
那么通常的稳定映射紧化
$\mathcal M^{\mathrm{main}}_{0,\ell}(L;\beta)$
未必紧。
下面细究这一现象。

首先，有边黎曼曲面的模空间之定义
包含亏格与边界分支数。
在光滑稳定映射的模空间的紧化中，
新添元素的定义域必须与邻近光滑元素
具有相同亏格和相同边界分支数。
本书只考虑亏格为零、
边界分支数为一的情形。

取稳定映射
$(\Sigma,w)=((|\Sigma|,\vec z,\vec z^+),w)
\in\mathcal M^{\mathrm{main}}_{k,\ell}(L;\beta)$；
为简便起见，假设 $k\ne0$。
利用保序单射
$\varrho^+\colon
\{1,\ldots,\ell'\}\to\{1,\ldots,\ell\}$，
忘掉所有边界标记点及一部分内部标记点，
并试图据此定义新的定义域，
得到亏格为零的有节点有边黎曼曲面
$\Sigma'=(|\Sigma|,\varnothing,\vec z^{+\prime})$。
现在考察
$(\Sigma',w)$ 的不稳定分支 $D$。
除第~\ref{sec:7-1} 节已经处理的
(7.1.41.1)--(7.1.41.3) 外，
还须单独处理：

\par\noindent\textup{(7.4.1)}\label{eq:7-4-1}\quad
只有一个内部奇点、
没有其他奇点或标记点的圆盘。

这种情形下不收缩 $D$。
按通常方式收缩其他类型的不稳定分支后，
所得对象或者属于
$\mathcal M^{\mathrm{main}}_{0,\ell'}(L;\beta)$，
或者有一个 (7.4.1) 型分支，
其上附着球面气泡。
于是得到遗忘映射
\begin{equation}
\begin{aligned}
\operatorname{Forget}_{\varnothing,\varrho^+}\colon&
\mathcal M^{\mathrm{main}}_{k,\ell}(L;\beta)\\
&\longrightarrow
\mathcal M^{\mathrm{main}}_{0,\ell'}(L;\beta)
\ \cup\
\bigcup_{\substack{\widetilde\beta:\\
 i_*(\widetilde\beta)=\beta}}
\bigl(\mathcal M_{\ell'+1}(M;\widetilde\beta)
\times_M L\bigr).
\end{aligned}
\tag{7.4.2}\label{eq:7-4-2}
\end{equation}
与第~\ref{subsec:3-8-3} 小节相同，
式~\eqref{eq:7-4-2} 的靶空间简记为
\[
\mathcal M^{\mathrm{main}}_{0,\ell'}(L;\beta)
\cup
\bigl(\mathcal M_{\ell'+1}
(M;i_*^{-1}(\beta))\times_M L\bigr).
\]
仿照定义~\ref{def:7-1-42}，
可在此并空间上定义拓扑；
仿照 \textbf{[FuOn99II]} 第~11、10 节，
可证它是紧豪斯多夫空间。

现在为
\[
\mathcal M^{\mathrm{main}}_{0,\ell}(L;\beta)
\cup
\bigl(\mathcal M_{\ell+1}
(M;i_*^{-1}(\beta))\times_M L\bigr)
\]
赋予带边界的仓西结构。
以下为简化记号，置 $\ell=\ell'$。
第~\ref{sec:7-1} 节已详细讨论
$\mathcal M^{\mathrm{main}}_{0,\ell}(L;\beta)$
各点的仓西邻域；
这里构造
$\mathcal M_{\ell+1}(M;\widetilde\beta)
\times_M L$ 上一点的仓西邻域，
其中 $\widetilde\beta\in i_*^{-1}(\beta)$。

设
$(\Sigma,w)=((\Sigma,z_0^+,\ldots,z_\ell^+),w)
\in\mathcal M_{\ell+1}(M;\widetilde\beta)$
是亏格为零、
带 $\ell+1$ 个标记点的稳定映射。
为简便起见，
只讨论 $\Sigma$ 光滑
（即微分同胚于 $S^2$）、
恰有两个标记点 $z_0^+,z_1^+$ 的情形。
假设 $w(z_0^+)\in L$。
再假设有一族由 $u\in U$ 参数化的
$(\Sigma_u,w_u)\in\mathcal M_2(M;\widetilde\beta)$，
且这一族横截；
特别地，
$\dim U$ 等于
$\mathcal M_2(M;\widetilde\beta)$ 的虚维数。
又假设
$\operatorname{ev}_0\colon U\to M$，
$\operatorname{ev}_0(u)=w_u(z_0^+)$
横截于 $L$，
且 $(\Sigma_u,w_u)$ 的自同构群平凡。
一般情形还需障碍丛及有限群作用；
障碍丛可如其他成熟情形一样处理，
这里略去。

我们要构造由
$(U_{\operatorname{ev}_0}\times_M L)
\times[0,\epsilon)$
参数化、属于
\[
\mathcal M^{\mathrm{main}}_{0,1}(L;\beta)
\cup
\bigl(\mathcal M_2(M;\widetilde\beta)
\times_M L\bigr)
\]
的一族元素。
其中
$(U_{\operatorname{ev}_0}\times_M L)\times\{0\}$
与
$\mathcal M_2(M;\widetilde\beta)\times_M L$
中的一个开集相识别。

取
$(D^2,z_0,z^+)$，
其中 $z_0\in\partial D^2$，
$z^+=0\in D^2$。
对满足 $\operatorname{ev}_0(u)\in L$ 的 $u\in U$，
令
$\Sigma_u^+=
\Sigma_u{}_{z_0^+}\#_{z^+}D^2$，
并定义映射
$w_u^+\colon\Sigma_u^+\to M$：
\[
w_u^+(z)=
\begin{cases}
w_u(z),&z\in\Sigma_u,\\
w_u(z_0^+),&z\in D^2.
\end{cases}
\]
于是有映射\translatornote{源书将 $\mathfrak I$
的定义域写作 $U$；
但只有 $\operatorname{ev}_0(u)\in L$
的参数才能附着所述常值圆盘，
并且后续映射均以
$U_{\operatorname{ev}_0}\times_M L$ 为定义域。
见疑误 A-0556。}
\[
\mathfrak I\colon
U_{\operatorname{ev}_0}\times_M L
\longrightarrow
\mathcal M^{\mathrm{main}}_{1,1}(L;\beta),
\qquad
\mathfrak I(u,x)=(\Sigma_u^+,w_u^+).
\]

由于 $\Sigma_u^+$ 有内部奇点，
可以用单位分解光滑化奇点，
以 $\mathbb C$ 中直径为 $\epsilon$
的小圆盘 $D^2(\epsilon)$
参数化其在
$\mathcal M^{\mathrm{main}}_{1,1}(L;\beta)$
附近的对象；
参见 \textbf{[FuOn99II]} 第~12 节。
\translatornote{源书把 $D^2(\epsilon)$
称为 $\mathbb C^2$ 中的小圆盘；
但下一页用 $S^1$ 作商并得到区间，
故此处应为一复维的 $\mathbb C$。
见疑误 A-0555。}
由此有映射
\begin{equation}
\mathfrak I^{+,\mathrm{pre}}\colon
(U_{\operatorname{ev}_0}\times_M L)
\times D^2(\epsilon)
\longrightarrow
\operatorname{Map}_{1,1}(L;\beta).
\tag{7.4.3}\label{eq:7-4-3}
\end{equation}
这里
$\operatorname{Map}_{1,1}(L;\beta)$
是满足
$\mathcal M^{\mathrm{main}}_{1,1}(L;\beta)$
全部条件、唯独不要求映射伪全纯的
“稳定映射”集合。
再由隐函数定理得到
\begin{equation}
\mathfrak I^+\colon
(U_{\operatorname{ev}_0}\times_M L)
\times D^2(\epsilon)
\longrightarrow
\mathcal M^{\mathrm{main}}_{1,1}(L;\beta).
\tag{7.4.4}\label{eq:7-4-4}
\end{equation}
\par\medskip\noindent\begin{minipage}{\textwidth}
\centering
\setlength{\unitlength}{0.75mm}
\begin{picture}(78,68)
\linethickness{0.65pt}
{\color{gray!55}
\qbezier(12,18)(39,37)(66,18)}
\qbezier(12,18)(39,-2)(66,18)
\qbezier(39,18)(17,18)(17,43)
\qbezier(17,43)(17,69)(39,69)
\qbezier(39,69)(61,69)(61,43)
\qbezier(61,43)(61,18)(39,18)
\put(39,18){\circle*{1.5}}
\put(66,18){\circle*{1.5}}
\put(37,42){$\beta$}
\put(69,15){$z_0$}
\end{picture}
\par\smallskip
{\renewcommand{\thefigure}{7.4.1}\refstepcounter{figure}
\normalfont\bfseries 图~\thefigure\label{fig:7-4-1}}
\end{minipage}\par\medskip


\begin{lemma}{7.4.5}\label{lem:7-4-5}
$\mathfrak I^+$ 是开嵌入。
\end{lemma}
\begin{proof}
我们已假设
$\mathcal M^{\mathrm{main}}_{1,1}(L;\beta)$
横截于 $L$。
\translatornote{源书此处称圆盘模空间本身
“横截于 $L$”；
前页明确的假设则是
$\operatorname{ev}_0\colon U\to M$
横截于 $L$。
尚不能确定原句应如何改写，
故保留源书表述；见疑误 A-0557。}
此外，
$\mathcal M^{\mathrm{main}}_{1,1}(L;\beta)$
的元素没有对称性。
因此，将
$\mathcal M^{\mathrm{main}}_{1,1}(L;\beta)$
与 $\mathcal M_2(M;\widetilde\beta)$ 胶合，
所得空间微分同胚于相应纤维积
加上胶合参数，
结论随之成立。
详情参见 \textbf{[FuOn99II]} 第~13 节。
\end{proof}

接着在 $\mathfrak I^+$ 的像上
考虑遗忘映射
\[
\operatorname{Forget}\colon
\mathcal M^{\mathrm{main}}_{1,1}(L;\beta)
\longrightarrow
\mathcal M^{\mathrm{main}}_{0,1}(L;\beta)
\cup
\bigl(\mathcal M_2(M;i_*^{-1}(\beta))
\times_M L\bigr).
\]
由 $\mathfrak I^{+,\mathrm{pre}}$ 的构造，
容易看出下述引理。
令
\[
\operatorname{Forget}\colon
\operatorname{Map}_{1,1}(L;\beta)
\longrightarrow
\operatorname{Map}_{0,1}(L;\beta)
\cup
\bigl(\operatorname{Map}_2(M;i_*^{-1}(\beta))
\times_M L\bigr)
\]
为显然的遗忘映射，
并考虑 $D^2(\epsilon)$ 上的
$S^1$ 作用。
\begin{lemma}{7.4.6}\label{lem:7-4-6}
对 $\theta\in S^1$，
\[
\operatorname{Forget}
\bigl(\mathfrak I^{+,\mathrm{pre}}((u,x),\theta z)\bigr)
=
\operatorname{Forget}
\bigl(\mathfrak I^{+,\mathrm{pre}}((u,x),z)\bigr),
\]
其中
$(u,x)\in U_{\operatorname{ev}_0}\times_M L$，
$z\in D^2(\epsilon)$。
\end{lemma}

再看从 $\mathfrak I^{+,\mathrm{pre}}$
经隐函数定理得到 $\mathfrak I^+$ 的过程。
该构造显然在 $S^1$ 作用下不变，
因此
\[
\operatorname{Forget}
\bigl(\mathfrak I^+((u,x),\theta z)\bigr)
=
\operatorname{Forget}
\bigl(\mathfrak I^+((u,x),z)\bigr).
\]
于是
$\mathfrak I^+$ 的像在
$\operatorname{Forget}$ 下的像
微分同胚于
\[
(U_{\operatorname{ev}_0}\times_M L)
\times D^2(\epsilon)/S^1
\cong
(U_{\operatorname{ev}_0}\times_M L)
\times[0,\epsilon).
\]
所需仓西邻域由此构造完成。
命题~\ref{prop:3-8-27} 证明的余下部分
与现已成熟的其他论证相同。\hfill$\square$
